\section{Formal verification architecture and trust boundary}
\label{sec:formalisation}

The fixed-ray and fixed-fibre core of this paper has been formalised in Lean~4,
an interactive theorem prover with a small proof-checking kernel
\parencite{deMouraUllrich2021}, using the Mathlib library of formalised
mathematics \parencite{MathlibCommunity2020}.  The formalisation is not a
translation of displayed formulas into numerical tests.  It defines the locked
six-state matrix, selected ports, perturbation, cofactor, word coefficients,
coefficient support and real convex hull, and then constructs proof terms for
the corresponding universal statements.

\subsection{End-to-end verified spine}

The organising dependency chain is recorded in \cref{tab:leanmap}.  Module
names and theorem identifiers are part of the release interface: they identify
the exact machine-checked statement rather than a nearby prose summary.

\begin{table}[ht]
\centering
\small
\begin{tabularx}{\textwidth}{@{}p{0.10\textwidth}p{0.34\textwidth}X@{}}
\toprule
Stage & Paper-level claim & Lean module and principal theorem\\
\midrule
M0--M1 & Locked model, signed involutions and all-depth baseline darkness &
\path{M3A.BaselineDarkness}; \path{baseline_moments_dark}\\
M2 & Exact unreduced selected cofactor numerator &
\path{M3A.SelectedNumerator}; \path{selectedNumerator_formula}\\
M3 & Pointwise decomposition into the two dark planes and persistent certificates &
\path{M3A.DarkPlanes}; \path{ell_quad_zero_iff_planes}\\
M4 & Locked endpoint-access support and signed boundary power sum &
\path{M3A.EndpointWords}; \path{sixState_support_bound},
\path{sixState_boundary_ray}\\
M5 & Exact first amplitude coefficients and complete first-order protection on $L=0$ &
\path{M3A.RayCoefficients};
\path{amplitudeCoeff_one_eq_zero_of_linear_zero}\\
M6 & Exhaustive linear/quadratic/dark fixed-ray classification &
\path{M3A.FixedRayTrichotomy}; \path{fixedRay_trichotomy},
\path{classifyRay_spec}\\
M7 & Actual-support Newton orthants and unique positive-weight vertices &
\path{M3A.FixedFibreNewton}; \path{fixedFibre_newton}\\
M8 & Identification with the selected matrix-exponential amplitude &
\path{M3A.AnalyticBridge};
\path{selectedRayAmplitude_eq_tsum_coefficients}\\
\bottomrule
\end{tabularx}
\caption{The theorem-to-Lean dependency map.  Each row depends on the rows
above it; M8 closes the gap between finite word algebra and the analytic
amplitude used in the paper.}
\label{tab:leanmap}
\end{table}

The resulting chain is genuinely end to end.  At each time depth $q$, M8
proves the finite identity
\[
M_q(\e v)=\sum_{p=0}^{q}\e^p M_{p,q}(v),
\]
where $M_{p,q}$ is the word coefficient defined in M4.  It then proves global
summability in $t$ and the equality
\begin{equation}\label{eq:leanbridge}
K(\e v,t)=
\sum_{q=0}^{\infty}\sum_{p=0}^{q}
A_{p,q}(v)\e^pt^q,
\qquad
A_{p,q}(v)=\frac{(-i)^q}{q!}M_{p,q}(v).
\end{equation}
Thus the support classified in M6 and convexified in M7 is the coefficient
support of the actual matrix exponential, not an unconnected formal surrogate.

\subsection{What the status does and does not mean}

The release pins Lean and Mathlib at version 4.32.1.  A clean build checks the
root import through \path{M3A.AnalyticBridge}.  Axiom inspection of every
principal theorem reports only \path{propext}, \path{Classical.choice} and
\path{Quot.sound}, the standard foundational axioms used by Lean and Mathlib.
The project contains no \path{sorry}, \path{admit}, user-declared axiom or
unsafe theorem dependency.  Exact Python and computer-algebra scripts remain
valuable independent regression checks, but they are not premises of the Lean
proofs.

The formal boundary is equally important.  M0--M8 cover the locked scalar-ray
specialisation through the fixed-fibre Newton theorem and its analytic bridge.
They do not yet formalise the fully general rectangular endpoint theorem, the
ideal equality in a multivariate polynomial ring, monomial pullback, moving
analytic arcs, the universal four-coordinate atlas, the Gr\"obner/Rees
structures or the fixed-weight Noetherian theorem.  Those later results retain
their conventional proofs and exact-computation certificates.  The transition
is marked explicitly after \cref{sec:fixednewton}, and the detailed claim map is
given in \cref{app:reprochecks}.
